\documentclass[preview]{standalone}

\usepackage{amsmath}
\usepackage{amssymb}
\usepackage{stellar}
\usepackage{definitions}
\usepackage{tikz}

\usetikzlibrary{cd}

\begin{document}

\id{quotient-spaces-isomorphism-theorems}
\genpage

\section{Canonical projection}

\begin{snippetproposition}{quotient-vector-space-canonical-projection}{Canonical projection onto a quotient space}
    Let \(V\) be a \vectorspace over a \field \(\mathbb{F}\), and let
    \(W\subseteq V\) be a linear subspace. The \function
    \[
        \pi\colon V\fromto V/W,
        \qquad
        \pi(v)=v+W,
    \]
    is a \surjective \lineartransformation with
    \(\ltkernel\pi=W\).
\end{snippetproposition}

\begin{snippetproof}{quotient-vector-space-canonical-projection-proof}{quotient-vector-space-canonical-projection}{Canonical projection onto a quotient space}
    For \(v_1,v_2\in V\) and \(\lambda\in\mathbb{F}\),
    \begin{align*}
        \pi(v_1+v_2)
        &=(v_1+v_2)+W
        =(v_1+W)+(v_2+W), \\
        \pi(\lambda v_1)
        &=\lambda v_1+W
        =\lambda(v_1+W).
    \end{align*}
    Thus \(\pi\) is linear. Every coset \(v+W\) equals \(\pi(v)\), so
    \(\pi\) is \surjective. Finally,
    \[
        \pi(v)=0_{V/W}
        \ifandonlyif v+W=W
        \ifandonlyif v\in W,
    \]
    whence \(\ltkernel\pi=W\).
\end{snippetproof}

\begin{snippetproposition}{linear-map-induced-on-kernel-quotient}{Linear map induced on the quotient by the kernel}
    Let \(V_1,V_2\) be \vectorspace[vector spaces] over a \field
    \(\mathbb{F}\), and let
    \(\varphi\colon V_1\fromto V_2\) be a \lineartransformation. The
    \function
    \[
        \overline{\varphi}\colon V_1/\ltkernel\varphi\fromto V_2,
        \qquad
        \overline{\varphi}(v+\ltkernel\varphi)=\varphi(v),
    \]
    is a well-defined \lineartransformation.
\end{snippetproposition}

\begin{snippetproof}{linear-map-induced-on-kernel-quotient-proof}{linear-map-induced-on-kernel-quotient}{Linear map induced on the quotient by the kernel}
    Suppose that \(v'=v+w\), where \(w\in\ltkernel\varphi\). Then
    \[
        \varphi(v')=\varphi(v)+\varphi(w)=\varphi(v),
    \]
    so \(\overline{\varphi}\) is well-defined. For \(v_1,v_2\in V_1\) and
    \(\lambda\in\mathbb{F}\),
    \begin{align*}
        \overline{\varphi}((v_1+\ltkernel\varphi)+(v_2+\ltkernel\varphi))
        &=\varphi(v_1+v_2) \\
        &=\overline{\varphi}(v_1+\ltkernel\varphi)
          +\overline{\varphi}(v_2+\ltkernel\varphi), \\
        \overline{\varphi}(\lambda(v_1+\ltkernel\varphi))
        &=\varphi(\lambda v_1)
        =\lambda\overline{\varphi}(v_1+\ltkernel\varphi).
    \end{align*}
\end{snippetproof}

\begin{snippet}{linear-map-kernel-quotient-factorization}
    Let \(V_1,V_2\) be \vectorspace[vector spaces] over a \field
    \(\mathbb{F}\), let
    \(\varphi\colon V_1\fromto V_2\) be a \lineartransformation, and let
    \(\pi\colon V_1\fromto V_1/\ltkernel\varphi\) be the canonical
    projection. Then \(\varphi=\overline{\varphi}\circ\pi\). Equivalently,
    the following diagram commutes:
    \[
        \begin{tikzcd}
            V_1 \arrow[r, "\varphi"] \arrow[d, "\pi"'] & V_2 \\
            V_1/\ltkernel\varphi
                \arrow[ur, "\overline{\varphi}"']
        \end{tikzcd}
    \]
    Indeed,
    \[
        (\overline{\varphi}\circ\pi)(v)
        =\overline{\varphi}(v+\ltkernel\varphi)
        =\varphi(v)
    \]
    for every \(v\in V_1\).
\end{snippet}

\section{Isomorphism theorems}

\begin{snippettheorem}{vector-space-first-isomorphism-theorem}{First isomorphism theorem for vector spaces}
    Let \(V_1,V_2\) be \vectorspace[vector spaces] over a \field
    \(\mathbb{F}\), and let
    \(\varphi\colon V_1\fromto V_2\) be a \lineartransformation. Then
    \[
        V_1/\ltkernel\varphi\cong\ltimage\varphi.
    \]
    More precisely, the \function
    \[
        \overline{\varphi}\colon
        V_1/\ltkernel\varphi\fromto\ltimage\varphi,
        \qquad
        \overline{\varphi}(v+\ltkernel\varphi)=\varphi(v),
    \]
    is a \linearisomorphism.
\end{snippettheorem}

\begin{snippetproof}{vector-space-first-isomorphism-theorem-proof}{vector-space-first-isomorphism-theorem}{First isomorphism theorem for vector spaces}
    The induced \function \(\overline{\varphi}\) is linear and well-defined.
    It is \surjective onto \(\ltimage\varphi\) by definition. Moreover,
    \begin{align*}
        \ltkernel\overline{\varphi}
        &=\{v+\ltkernel\varphi
            \suchthat\varphi(v)=0_{V_2}\} \\
        &=\{0_{V_1}+\ltkernel\varphi\}.
    \end{align*}
    Hence \(\overline{\varphi}\) is \injective and therefore a
    \linearisomorphism.
\end{snippetproof}

\begin{snippetexample}{direct-sum-quotient-isomorphism-example}{Quotients of a direct sum}
    Let \(V\) be a \vectorspace over a \field \(\mathbb{F}\), and suppose
    that \(V=W_1\oplus W_2\). Every \(v\in V\) has a unique expression
    \(v=w_1+w_2\), with \(w_i\in W_i\). The projection
    \[
        p_1\colon V\fromto W_1,
        \qquad
        p_1(w_1+w_2)=w_1,
    \]
    has kernel \(W_2\), while the analogous projection \(p_2\) has kernel
    \(W_1\). The first isomorphism theorem therefore gives
    \[
        V/W_2\cong W_1,
        \qquad
        V/W_1\cong W_2.
    \]
    The corresponding factorizations are
    \[
        \begin{tikzcd}[column sep=large]
            V \arrow[r, "p_1"] \arrow[d, "\pi_1"'] & W_1 \\
            V/W_2 \arrow[ur, "\overline{p_1}"']
        \end{tikzcd}
        \qquad
        \begin{tikzcd}[column sep=large]
            V \arrow[r, "p_2"] \arrow[d, "\pi_2"'] & W_2 \\
            V/W_1 \arrow[ur, "\overline{p_2}"']
        \end{tikzcd}
    \]
\end{snippetexample}

\begin{snippettheorem}{vector-space-second-isomorphism-theorem}{Second isomorphism theorem for vector spaces}
    Let \(V\) be a \vectorspace over a \field \(\mathbb{F}\), and let
    \(W,U\subseteq V\) be linear subspaces. Then
    \[
        W/(W\intersection U)\cong(W+U)/U.
    \]
\end{snippettheorem}

\begin{snippetproof}{vector-space-second-isomorphism-theorem-proof}{vector-space-second-isomorphism-theorem}{Second isomorphism theorem for vector spaces}
    Consider the \lineartransformation
    \[
        \varphi\colon W\fromto(W+U)/U,
        \qquad
        \varphi(w)=w+U.
    \]
    Every element of \((W+U)/U\) has the form
    \((w+u)+U=w+U\), where \(w\in W\) and \(u\in U\). Thus
    \(\varphi\) is \surjective. Furthermore,
    \begin{align*}
        \ltkernel\varphi
        &=\{w\in W\suchthat w+U=U\} \\
        &=\{w\in W\suchthat w\in U\}
        =W\intersection U.
    \end{align*}
    The first isomorphism theorem now yields
    \[
        W/(W\intersection U)
        =W/\ltkernel\varphi
        \cong\ltimage\varphi
        =(W+U)/U.
    \]
\end{snippetproof}

\begin{snippet}{vector-space-second-isomorphism-theorem-inclusion-note}
    Let \(U\) and \(W\) be linear subspaces of a \vectorspace \(V\). The
    second isomorphism theorem does not require \(U\subseteq W\) or
    \(W\subseteq U\). In particular, it compares the two subspaces even when
    neither quotient \(W/U\) nor \(U/W\) is defined.
\end{snippet}

\begin{snippettheorem}{vector-space-third-isomorphism-theorem}{Third isomorphism theorem for vector spaces}
    Let \(V\) be a \vectorspace over a \field \(\mathbb{F}\), and let
    \(U\subseteq W\subseteq V\) be linear subspaces. Then \(W/U\) is a
    linear subspace of \(V/U\), and
    \[
        (V/U)/(W/U)\cong V/W.
    \]
\end{snippettheorem}

\begin{snippetproof}{vector-space-third-isomorphism-theorem-proof}{vector-space-third-isomorphism-theorem}{Third isomorphism theorem for vector spaces}
    Since \(U\subseteq W\subseteq V\), the operations of \(W/U\) are the
    restrictions of those of \(V/U\), so \(W/U\) is a linear subspace of
    \(V/U\). Define
    \[
        \overline{\varphi}\colon V/U\fromto V/W,
        \qquad
        \overline{\varphi}(v+U)=v+W.
    \]
    If \(v'+U=v+U\), then \(v'-v\in U\subseteq W\), whence
    \(v'+W=v+W\). Thus \(\overline{\varphi}\) is well-defined. Its
    linearity follows directly from the quotient operations, and it is
    \surjective because every \(v+W\) is the image of \(v+U\). Finally,
    \begin{align*}
        \ltkernel\overline{\varphi}
        &=\{v+U\in V/U\suchthat v+W=W\} \\
        &=\{v+U\in V/U\suchthat v\in W\}
        =W/U.
    \end{align*}
    The first isomorphism theorem gives
    \[
        (V/U)/(W/U)
        =(V/U)/\ltkernel\overline{\varphi}
        \cong\ltimage\overline{\varphi}
        =V/W.
    \]
\end{snippetproof}

\end{document}
